% Фрагмент для вставки в отчёт: % Фрагмент для вставки в отчёт: % Фрагмент для вставки в отчёт: % Фрагмент для вставки в отчёт: \input{chapters/parsing_cky_earley}
% Источники: report.pdf (разд.~2.3.1), Jurafsky \& Martin, 3rd draft, гл.~13;
% Earley — Earley (1970); в 3-м черновике Jurafsky (2026) — краткое упоминание в гл.~13.

\subsection{Алгоритм Кока--Янгера--Касами (CKY)}

\subsubsection{Что можно сохранить из текущего описания (\texttt{report.pdf}, разд.~2.3.1)}

Следующие положения согласуются с изложением в~\cite{jurafsky2026speech}, гл.~13, и могут быть оставлены (после небольшой правки формулировок):
\begin{itemize}
	\item алгоритм строит \emph{дерево составляющих} (конституентный разбор) для предложений, порождаемых контекстно-свободной грамматикой;
	\item перед разбором грамматику приводят к \emph{нормальной форме Хомского} (НФХ): правила вида $A \to BC$ или $A \to w$~\cite{jurafsky2026speech};
	\item используется \emph{динамическое программирование} на верхнетреугольной таблице размера $(n{+}1)\times(n{+}1)$ для предложения из $n$ слов;
	\item ответ разбора (все допустимые корневые конституенты) содержится в ячейке $[0,n]$;
	\item для \emph{восстановления дерева} целесообразно хранить указатели на подынтервалы $(i,k)$ и $(k,j)$, породившие данную запись;
	\item \emph{временная сложность} распознавания --- $O(n^3)$ при фиксированной грамматике (точнее $O(n^3|G|)$, если учитывать размер грамматики).
\end{itemize}

\subsubsection{Замечания и ошибки относительно~\cite{jurafsky2026speech}, гл.~13}

\begin{enumerate}
	\item \textbf{Название алгоритма.}
	В отчёте: «Кок--Янгер--Касами--Шварц--Зиппель».
	В гл.~13 Jurafsky --- классический \textbf{CKY / CYK} (Cocke--Kasami--Younger)~\cite{younger1967recognition}.
	Расширение Schwartz--Zippel --- отдельное современное направление~\cite{makarov2022cyk}; для базового парсера его включать в название не следует.

	\item \textbf{Рекуррентное соотношение (существенная ошибка).}
	В отчёте:
	\[
		dp_{i,j} = dp_{i,j} \cup \bigl(dp_{i,k} \cup dp_{k,j}\bigr), \quad i<k<j,
	\]
	что не задаёт порождение новых нетерминалов.
	По~\cite{jurafsky2026speech}, \S13.2.2--13.2.3, ячейка $[i,j]$ --- это \emph{множество нетерминалов}, покрывающих слова с позиции $i$ по $j$.
	Для каждого разбиения $i<k<j$ и каждого правила $A \to B\,C$ из НФХ, если $B \in \mathrm{cell}[i,k]$ и $C \in \mathrm{cell}[k,j]$, добавляют $A$ в $\mathrm{cell}[i,j]$.

	\item \textbf{Семантика ячеек.}
	В отчёте смешаны «тип объединения» и объединение множеств без ссылки на правила грамматики.
	В Jurafsky объединяются только те нетерминалы, которые \emph{разрешены грамматикой} при бинарном разбиении интервала.

	\item \textbf{Порядок заполнения таблицы.}
	Формулировка «слева направо, снизу вверх» неточна.
	В гл.~13: внешний цикл --- по \emph{столбцам} (длине интервала) слева направо, внутри столбца --- по \emph{строкам} снизу вверх; для ячейки $[i,j]$ перебирают все точки разбиения $k$ ($i<k<j$).

	\item \textbf{Индексация.}
	У Jurafsky индексы $0,\ldots,n$ --- «столбики» \emph{между} словами; слова дают заполнение диагонали вида $[i,i{+}1]$.
	В отчёте $i,j = \overline{0,n}$ при $i<j$ без явного различения ячеек длины~1 и произвольного span --- это стоит уточнить.

	\item \textbf{Распознавание и разбор.}
	Базовый CKY --- \emph{распознаватель}: успех, если $S \in \mathrm{cell}[0,n]$ (\S13.2.3).
	Фраза «варианты разбора в $dp_{0,n}$» верна только если в ячейках хранят \emph{все} нетерминалы и (для полного парсера) обратные указатели; иначе это лишь множество возможных корней.

	\item \textbf{Дисамбигуация.}
	«Вручную или предобученным классификатором» --- не часть классического CKY в \S13.2.
	Выбор лучшего разбора Jurafsky выносит в \S13.3+ (нейросетевой span-based parsing, вероятностный CKY); для диплома формулировку лучше отделить от описания ядра алгоритма.

	\item \textbf{Ссылки на доказательство сложности.}
	Доказательство $O(n^3)$ для CYK --- Younger (1967), Kasami (1965); в~\cite{makarov2022cyk} --- обзор родственных алгоритмов, а не классическое доказательство.
	Теорему о приведении к НФХ логичнее ссылать на гл.~12~\cite{jurafsky2026speech}, а не только на Makarov.

	\item \textbf{Опечатки в соседнем тексте.}
	«завсисмостей» (табл.~сравнения), «было было» (восстановление дерева).
\end{enumerate}

\subsubsection{Исправленное описание (по~\cite{jurafsky2026speech}, гл.~13)}

Пусть вход --- слова $w_1,\ldots,w_n$, грамматика $G=(N,\Sigma,P,S)$ в НФХ.
Таблица $\mathrm{cell}[i,j]$ ($0 \le i < j \le n$) --- множество нетерминалов, порождающих $w_{i+1}\ldots w_j$.

\paragraph{Инициализация.}
Для каждого $i=0,\ldots,n{-}1$:
$\mathrm{cell}[i,i{+}1] \gets \{A \mid (A \to w_{i+1}) \in P\}$ (и при необходимости частеречные категории из лексикона).

\paragraph{Заполнение.}
Для $j = 2,\ldots,n$ (длина интервала), для $i = j{-}2,\ldots,0$ (снизу вверх в столбце $j$), для $k = i{+}1,\ldots,j{-}1$:
\[
	\mathrm{cell}[i,j] \gets \mathrm{cell}[i,j] \cup
	\{A \mid A \to B\,C \in P,\; B \in \mathrm{cell}[i,k],\; C \in \mathrm{cell}[k,j]\}.
\]

\paragraph{Распознавание.}
Предложение принимается грамматикой тогда и только тогда, когда $S \in \mathrm{cell}[0,n]$.

\paragraph{Построение дерева.}
При добавлении $A$ из правила $A \to B\,C$ с разбиением $k$ сохраняют указатель $(B,C,k)$.
Восстановление --- рекурсивный обход от $S \in \mathrm{cell}[0,n]$.

\paragraph{Сложность.}
При $|G|$ правил и длине $n$ --- $O(n^3|G|)$; при фиксированной грамматике --- $O(n^3)$.

%--------------------------------------------------------------------
\subsection{Алгоритм Эрли (Earley)}

% В 3-м черновике Jurafsky (2026) алгоритм Эрли упоминается в гл.~13
% как родственный подход; детали ниже по~\cite{earley1970efficient}.

Алгоритм Эрли~\cite{earley1970efficient} --- алгоритм \emph{динамического программирования} для разбора произвольных КС-грамматик (НФХ \emph{не требуется}).
Он совмещает нисходящие ожидания и восходящее подтверждение конституентов; в отличие от CKY, заполнение идёт \emph{слева направо} по позициям входа.

\subsubsection{Структуры данных}

Для предложения длины $n$ строят \textbf{чарт} (chart) --- массив из $n{+}1$ множеств состояний $\mathrm{chart}[0],\ldots,\mathrm{chart}[n]$.
Позиция $i$ соответствует \emph{зазору} перед $i$-м словом (между словами).

\textbf{Состояние} (dotted rule) записывают как тройку
\[
	(A \to \alpha \bullet \beta;\; [i,j]),
\]
где $A \to \alpha\beta$ --- правило грамматики, $\bullet$ --- позиция в правой части, $[i,j]$ --- интервал входа, который уже покрыт префиксом $\alpha$ (от $i$ до текущей позиции обработки).

Стартовое состояние: $(\gamma \to \bullet S;\,[0,0])$, где $\gamma$ --- фиктивный начальный символ.

\subsubsection{Операторы}

Для каждой позиции $i$ обрабатывают состояния в $\mathrm{chart}[i]$ в порядке добавления (без удаления и без возврата назад).

\begin{description}
	\item[\textsc{Predictor}] применяют, если справа от $\bullet$ стоит нетерминал $B$ (не часть речи).
	Для каждого правила $B \to \gamma$ добавляют $(B \to \bullet \gamma;\,[j,j])$ в $\mathrm{chart}[j]$, где $j$ --- конец уже разобранного префикса состояния.

	\item[\textsc{Scanner}] применяют, если справа от $\bullet$ --- терминал/часть речи $B$ и $w_i$ согласуется с $B$.
	Создают $(B \to w_i \bullet;\,[i,i{+}1])$ в $\mathrm{chart}[i{+}1]$.

	\item[\textsc{Completer}] применяют, если $\bullet$ в конце правила: $(B \to \gamma \bullet;\,[j,k])$ завершено.
	Для каждого $(A \to \alpha \bullet B \beta;\,[i,j])$ в $\mathrm{chart}[j]$ добавляют $(A \to \alpha B \bullet \beta;\,[i,k])$ в $\mathrm{chart}[k]$.
\end{description}

Дубликаты состояний не добавляют (\textsc{Enqueue} только при отсутствии состояния в множестве).

\subsubsection{Результат}

Разбор успешен, если в $\mathrm{chart}[n]$ есть состояние $(\gamma \to S \bullet;\,[0,n])$ (или эквивалентно завершённое $S$, покрывающее всё предложение).
Чарт неявно кодирует \emph{все} синтаксические деревья; поддеревья хранятся один раз и переиспользуются.

Для извлечения конкретного дерева модифицируют \textsc{Completer}, сохраняя ссылки на породившие состояния (аналогично обратным указателям в CKY).

\subsubsection{Сложность и сравнение с CKY}

\begin{itemize}
	\item общий случай: $O(n^3)$ по длине $n$;
	\item для неоднозначных грамматик на практике часто ближе к $O(n^2)$;
	\item для некоторых классов грамматик (например, $LR(k)$) --- линейное время;
	\item в отличие от CKY, не нужна НФХ; зато реализация и анализ состояний сложнее.
\end{itemize}

В~\cite{jurafsky2026speech}, гл.~13, CKY излагается как основной табличный bottom-up алгоритм; алгоритм Эрли и chart parsing упоминаются как родственные подходы на основе того же принципа мемоизации поддеревьев.

% Источники: report.pdf (разд.~2.3.1), Jurafsky \& Martin, 3rd draft, гл.~13;
% Earley — Earley (1970); в 3-м черновике Jurafsky (2026) — краткое упоминание в гл.~13.

\subsection{Алгоритм Кока--Янгера--Касами (CKY)}

\subsubsection{Что можно сохранить из текущего описания (\texttt{report.pdf}, разд.~2.3.1)}

Следующие положения согласуются с изложением в~\cite{jurafsky2026speech}, гл.~13, и могут быть оставлены (после небольшой правки формулировок):
\begin{itemize}
	\item алгоритм строит \emph{дерево составляющих} (конституентный разбор) для предложений, порождаемых контекстно-свободной грамматикой;
	\item перед разбором грамматику приводят к \emph{нормальной форме Хомского} (НФХ): правила вида $A \to BC$ или $A \to w$~\cite{jurafsky2026speech};
	\item используется \emph{динамическое программирование} на верхнетреугольной таблице размера $(n{+}1)\times(n{+}1)$ для предложения из $n$ слов;
	\item ответ разбора (все допустимые корневые конституенты) содержится в ячейке $[0,n]$;
	\item для \emph{восстановления дерева} целесообразно хранить указатели на подынтервалы $(i,k)$ и $(k,j)$, породившие данную запись;
	\item \emph{временная сложность} распознавания --- $O(n^3)$ при фиксированной грамматике (точнее $O(n^3|G|)$, если учитывать размер грамматики).
\end{itemize}

\subsubsection{Замечания и ошибки относительно~\cite{jurafsky2026speech}, гл.~13}

\begin{enumerate}
	\item \textbf{Название алгоритма.}
	В отчёте: «Кок--Янгер--Касами--Шварц--Зиппель».
	В гл.~13 Jurafsky --- классический \textbf{CKY / CYK} (Cocke--Kasami--Younger)~\cite{younger1967recognition}.
	Расширение Schwartz--Zippel --- отдельное современное направление~\cite{makarov2022cyk}; для базового парсера его включать в название не следует.

	\item \textbf{Рекуррентное соотношение (существенная ошибка).}
	В отчёте:
	\[
		dp_{i,j} = dp_{i,j} \cup \bigl(dp_{i,k} \cup dp_{k,j}\bigr), \quad i<k<j,
	\]
	что не задаёт порождение новых нетерминалов.
	По~\cite{jurafsky2026speech}, \S13.2.2--13.2.3, ячейка $[i,j]$ --- это \emph{множество нетерминалов}, покрывающих слова с позиции $i$ по $j$.
	Для каждого разбиения $i<k<j$ и каждого правила $A \to B\,C$ из НФХ, если $B \in \mathrm{cell}[i,k]$ и $C \in \mathrm{cell}[k,j]$, добавляют $A$ в $\mathrm{cell}[i,j]$.

	\item \textbf{Семантика ячеек.}
	В отчёте смешаны «тип объединения» и объединение множеств без ссылки на правила грамматики.
	В Jurafsky объединяются только те нетерминалы, которые \emph{разрешены грамматикой} при бинарном разбиении интервала.

	\item \textbf{Порядок заполнения таблицы.}
	Формулировка «слева направо, снизу вверх» неточна.
	В гл.~13: внешний цикл --- по \emph{столбцам} (длине интервала) слева направо, внутри столбца --- по \emph{строкам} снизу вверх; для ячейки $[i,j]$ перебирают все точки разбиения $k$ ($i<k<j$).

	\item \textbf{Индексация.}
	У Jurafsky индексы $0,\ldots,n$ --- «столбики» \emph{между} словами; слова дают заполнение диагонали вида $[i,i{+}1]$.
	В отчёте $i,j = \overline{0,n}$ при $i<j$ без явного различения ячеек длины~1 и произвольного span --- это стоит уточнить.

	\item \textbf{Распознавание и разбор.}
	Базовый CKY --- \emph{распознаватель}: успех, если $S \in \mathrm{cell}[0,n]$ (\S13.2.3).
	Фраза «варианты разбора в $dp_{0,n}$» верна только если в ячейках хранят \emph{все} нетерминалы и (для полного парсера) обратные указатели; иначе это лишь множество возможных корней.

	\item \textbf{Дисамбигуация.}
	«Вручную или предобученным классификатором» --- не часть классического CKY в \S13.2.
	Выбор лучшего разбора Jurafsky выносит в \S13.3+ (нейросетевой span-based parsing, вероятностный CKY); для диплома формулировку лучше отделить от описания ядра алгоритма.

	\item \textbf{Ссылки на доказательство сложности.}
	Доказательство $O(n^3)$ для CYK --- Younger (1967), Kasami (1965); в~\cite{makarov2022cyk} --- обзор родственных алгоритмов, а не классическое доказательство.
	Теорему о приведении к НФХ логичнее ссылать на гл.~12~\cite{jurafsky2026speech}, а не только на Makarov.

	\item \textbf{Опечатки в соседнем тексте.}
	«завсисмостей» (табл.~сравнения), «было было» (восстановление дерева).
\end{enumerate}

\subsubsection{Исправленное описание (по~\cite{jurafsky2026speech}, гл.~13)}

Пусть вход --- слова $w_1,\ldots,w_n$, грамматика $G=(N,\Sigma,P,S)$ в НФХ.
Таблица $\mathrm{cell}[i,j]$ ($0 \le i < j \le n$) --- множество нетерминалов, порождающих $w_{i+1}\ldots w_j$.

\paragraph{Инициализация.}
Для каждого $i=0,\ldots,n{-}1$:
$\mathrm{cell}[i,i{+}1] \gets \{A \mid (A \to w_{i+1}) \in P\}$ (и при необходимости частеречные категории из лексикона).

\paragraph{Заполнение.}
Для $j = 2,\ldots,n$ (длина интервала), для $i = j{-}2,\ldots,0$ (снизу вверх в столбце $j$), для $k = i{+}1,\ldots,j{-}1$:
\[
	\mathrm{cell}[i,j] \gets \mathrm{cell}[i,j] \cup
	\{A \mid A \to B\,C \in P,\; B \in \mathrm{cell}[i,k],\; C \in \mathrm{cell}[k,j]\}.
\]

\paragraph{Распознавание.}
Предложение принимается грамматикой тогда и только тогда, когда $S \in \mathrm{cell}[0,n]$.

\paragraph{Построение дерева.}
При добавлении $A$ из правила $A \to B\,C$ с разбиением $k$ сохраняют указатель $(B,C,k)$.
Восстановление --- рекурсивный обход от $S \in \mathrm{cell}[0,n]$.

\paragraph{Сложность.}
При $|G|$ правил и длине $n$ --- $O(n^3|G|)$; при фиксированной грамматике --- $O(n^3)$.

%--------------------------------------------------------------------
\subsection{Алгоритм Эрли (Earley)}

% В 3-м черновике Jurafsky (2026) алгоритм Эрли упоминается в гл.~13
% как родственный подход; детали ниже по~\cite{earley1970efficient}.

Алгоритм Эрли~\cite{earley1970efficient} --- алгоритм \emph{динамического программирования} для разбора произвольных КС-грамматик (НФХ \emph{не требуется}).
Он совмещает нисходящие ожидания и восходящее подтверждение конституентов; в отличие от CKY, заполнение идёт \emph{слева направо} по позициям входа.

\subsubsection{Структуры данных}

Для предложения длины $n$ строят \textbf{чарт} (chart) --- массив из $n{+}1$ множеств состояний $\mathrm{chart}[0],\ldots,\mathrm{chart}[n]$.
Позиция $i$ соответствует \emph{зазору} перед $i$-м словом (между словами).

\textbf{Состояние} (dotted rule) записывают как тройку
\[
	(A \to \alpha \bullet \beta;\; [i,j]),
\]
где $A \to \alpha\beta$ --- правило грамматики, $\bullet$ --- позиция в правой части, $[i,j]$ --- интервал входа, который уже покрыт префиксом $\alpha$ (от $i$ до текущей позиции обработки).

Стартовое состояние: $(\gamma \to \bullet S;\,[0,0])$, где $\gamma$ --- фиктивный начальный символ.

\subsubsection{Операторы}

Для каждой позиции $i$ обрабатывают состояния в $\mathrm{chart}[i]$ в порядке добавления (без удаления и без возврата назад).

\begin{description}
	\item[\textsc{Predictor}] применяют, если справа от $\bullet$ стоит нетерминал $B$ (не часть речи).
	Для каждого правила $B \to \gamma$ добавляют $(B \to \bullet \gamma;\,[j,j])$ в $\mathrm{chart}[j]$, где $j$ --- конец уже разобранного префикса состояния.

	\item[\textsc{Scanner}] применяют, если справа от $\bullet$ --- терминал/часть речи $B$ и $w_i$ согласуется с $B$.
	Создают $(B \to w_i \bullet;\,[i,i{+}1])$ в $\mathrm{chart}[i{+}1]$.

	\item[\textsc{Completer}] применяют, если $\bullet$ в конце правила: $(B \to \gamma \bullet;\,[j,k])$ завершено.
	Для каждого $(A \to \alpha \bullet B \beta;\,[i,j])$ в $\mathrm{chart}[j]$ добавляют $(A \to \alpha B \bullet \beta;\,[i,k])$ в $\mathrm{chart}[k]$.
\end{description}

Дубликаты состояний не добавляют (\textsc{Enqueue} только при отсутствии состояния в множестве).

\subsubsection{Результат}

Разбор успешен, если в $\mathrm{chart}[n]$ есть состояние $(\gamma \to S \bullet;\,[0,n])$ (или эквивалентно завершённое $S$, покрывающее всё предложение).
Чарт неявно кодирует \emph{все} синтаксические деревья; поддеревья хранятся один раз и переиспользуются.

Для извлечения конкретного дерева модифицируют \textsc{Completer}, сохраняя ссылки на породившие состояния (аналогично обратным указателям в CKY).

\subsubsection{Сложность и сравнение с CKY}

\begin{itemize}
	\item общий случай: $O(n^3)$ по длине $n$;
	\item для неоднозначных грамматик на практике часто ближе к $O(n^2)$;
	\item для некоторых классов грамматик (например, $LR(k)$) --- линейное время;
	\item в отличие от CKY, не нужна НФХ; зато реализация и анализ состояний сложнее.
\end{itemize}

В~\cite{jurafsky2026speech}, гл.~13, CKY излагается как основной табличный bottom-up алгоритм; алгоритм Эрли и chart parsing упоминаются как родственные подходы на основе того же принципа мемоизации поддеревьев.

% Источники: report.pdf (разд.~2.3.1), Jurafsky \& Martin, 3rd draft, гл.~13;
% Earley — Earley (1970); в 3-м черновике Jurafsky (2026) — краткое упоминание в гл.~13.

\subsection{Алгоритм Кока--Янгера--Касами (CKY)}

\subsubsection{Что можно сохранить из текущего описания (\texttt{report.pdf}, разд.~2.3.1)}

Следующие положения согласуются с изложением в~\cite{jurafsky2026speech}, гл.~13, и могут быть оставлены (после небольшой правки формулировок):
\begin{itemize}
	\item алгоритм строит \emph{дерево составляющих} (конституентный разбор) для предложений, порождаемых контекстно-свободной грамматикой;
	\item перед разбором грамматику приводят к \emph{нормальной форме Хомского} (НФХ): правила вида $A \to BC$ или $A \to w$~\cite{jurafsky2026speech};
	\item используется \emph{динамическое программирование} на верхнетреугольной таблице размера $(n{+}1)\times(n{+}1)$ для предложения из $n$ слов;
	\item ответ разбора (все допустимые корневые конституенты) содержится в ячейке $[0,n]$;
	\item для \emph{восстановления дерева} целесообразно хранить указатели на подынтервалы $(i,k)$ и $(k,j)$, породившие данную запись;
	\item \emph{временная сложность} распознавания --- $O(n^3)$ при фиксированной грамматике (точнее $O(n^3|G|)$, если учитывать размер грамматики).
\end{itemize}

\subsubsection{Замечания и ошибки относительно~\cite{jurafsky2026speech}, гл.~13}

\begin{enumerate}
	\item \textbf{Название алгоритма.}
	В отчёте: «Кок--Янгер--Касами--Шварц--Зиппель».
	В гл.~13 Jurafsky --- классический \textbf{CKY / CYK} (Cocke--Kasami--Younger)~\cite{younger1967recognition}.
	Расширение Schwartz--Zippel --- отдельное современное направление~\cite{makarov2022cyk}; для базового парсера его включать в название не следует.

	\item \textbf{Рекуррентное соотношение (существенная ошибка).}
	В отчёте:
	\[
		dp_{i,j} = dp_{i,j} \cup \bigl(dp_{i,k} \cup dp_{k,j}\bigr), \quad i<k<j,
	\]
	что не задаёт порождение новых нетерминалов.
	По~\cite{jurafsky2026speech}, \S13.2.2--13.2.3, ячейка $[i,j]$ --- это \emph{множество нетерминалов}, покрывающих слова с позиции $i$ по $j$.
	Для каждого разбиения $i<k<j$ и каждого правила $A \to B\,C$ из НФХ, если $B \in \mathrm{cell}[i,k]$ и $C \in \mathrm{cell}[k,j]$, добавляют $A$ в $\mathrm{cell}[i,j]$.

	\item \textbf{Семантика ячеек.}
	В отчёте смешаны «тип объединения» и объединение множеств без ссылки на правила грамматики.
	В Jurafsky объединяются только те нетерминалы, которые \emph{разрешены грамматикой} при бинарном разбиении интервала.

	\item \textbf{Порядок заполнения таблицы.}
	Формулировка «слева направо, снизу вверх» неточна.
	В гл.~13: внешний цикл --- по \emph{столбцам} (длине интервала) слева направо, внутри столбца --- по \emph{строкам} снизу вверх; для ячейки $[i,j]$ перебирают все точки разбиения $k$ ($i<k<j$).

	\item \textbf{Индексация.}
	У Jurafsky индексы $0,\ldots,n$ --- «столбики» \emph{между} словами; слова дают заполнение диагонали вида $[i,i{+}1]$.
	В отчёте $i,j = \overline{0,n}$ при $i<j$ без явного различения ячеек длины~1 и произвольного span --- это стоит уточнить.

	\item \textbf{Распознавание и разбор.}
	Базовый CKY --- \emph{распознаватель}: успех, если $S \in \mathrm{cell}[0,n]$ (\S13.2.3).
	Фраза «варианты разбора в $dp_{0,n}$» верна только если в ячейках хранят \emph{все} нетерминалы и (для полного парсера) обратные указатели; иначе это лишь множество возможных корней.

	\item \textbf{Дисамбигуация.}
	«Вручную или предобученным классификатором» --- не часть классического CKY в \S13.2.
	Выбор лучшего разбора Jurafsky выносит в \S13.3+ (нейросетевой span-based parsing, вероятностный CKY); для диплома формулировку лучше отделить от описания ядра алгоритма.

	\item \textbf{Ссылки на доказательство сложности.}
	Доказательство $O(n^3)$ для CYK --- Younger (1967), Kasami (1965); в~\cite{makarov2022cyk} --- обзор родственных алгоритмов, а не классическое доказательство.
	Теорему о приведении к НФХ логичнее ссылать на гл.~12~\cite{jurafsky2026speech}, а не только на Makarov.

	\item \textbf{Опечатки в соседнем тексте.}
	«завсисмостей» (табл.~сравнения), «было было» (восстановление дерева).
\end{enumerate}

\subsubsection{Исправленное описание (по~\cite{jurafsky2026speech}, гл.~13)}

Пусть вход --- слова $w_1,\ldots,w_n$, грамматика $G=(N,\Sigma,P,S)$ в НФХ.
Таблица $\mathrm{cell}[i,j]$ ($0 \le i < j \le n$) --- множество нетерминалов, порождающих $w_{i+1}\ldots w_j$.

\paragraph{Инициализация.}
Для каждого $i=0,\ldots,n{-}1$:
$\mathrm{cell}[i,i{+}1] \gets \{A \mid (A \to w_{i+1}) \in P\}$ (и при необходимости частеречные категории из лексикона).

\paragraph{Заполнение.}
Для $j = 2,\ldots,n$ (длина интервала), для $i = j{-}2,\ldots,0$ (снизу вверх в столбце $j$), для $k = i{+}1,\ldots,j{-}1$:
\[
	\mathrm{cell}[i,j] \gets \mathrm{cell}[i,j] \cup
	\{A \mid A \to B\,C \in P,\; B \in \mathrm{cell}[i,k],\; C \in \mathrm{cell}[k,j]\}.
\]

\paragraph{Распознавание.}
Предложение принимается грамматикой тогда и только тогда, когда $S \in \mathrm{cell}[0,n]$.

\paragraph{Построение дерева.}
При добавлении $A$ из правила $A \to B\,C$ с разбиением $k$ сохраняют указатель $(B,C,k)$.
Восстановление --- рекурсивный обход от $S \in \mathrm{cell}[0,n]$.

\paragraph{Сложность.}
При $|G|$ правил и длине $n$ --- $O(n^3|G|)$; при фиксированной грамматике --- $O(n^3)$.

%--------------------------------------------------------------------
\subsection{Алгоритм Эрли (Earley)}

% В 3-м черновике Jurafsky (2026) алгоритм Эрли упоминается в гл.~13
% как родственный подход; детали ниже по~\cite{earley1970efficient}.

Алгоритм Эрли~\cite{earley1970efficient} --- алгоритм \emph{динамического программирования} для разбора произвольных КС-грамматик (НФХ \emph{не требуется}).
Он совмещает нисходящие ожидания и восходящее подтверждение конституентов; в отличие от CKY, заполнение идёт \emph{слева направо} по позициям входа.

\subsubsection{Структуры данных}

Для предложения длины $n$ строят \textbf{чарт} (chart) --- массив из $n{+}1$ множеств состояний $\mathrm{chart}[0],\ldots,\mathrm{chart}[n]$.
Позиция $i$ соответствует \emph{зазору} перед $i$-м словом (между словами).

\textbf{Состояние} (dotted rule) записывают как тройку
\[
	(A \to \alpha \bullet \beta;\; [i,j]),
\]
где $A \to \alpha\beta$ --- правило грамматики, $\bullet$ --- позиция в правой части, $[i,j]$ --- интервал входа, который уже покрыт префиксом $\alpha$ (от $i$ до текущей позиции обработки).

Стартовое состояние: $(\gamma \to \bullet S;\,[0,0])$, где $\gamma$ --- фиктивный начальный символ.

\subsubsection{Операторы}

Для каждой позиции $i$ обрабатывают состояния в $\mathrm{chart}[i]$ в порядке добавления (без удаления и без возврата назад).

\begin{description}
	\item[\textsc{Predictor}] применяют, если справа от $\bullet$ стоит нетерминал $B$ (не часть речи).
	Для каждого правила $B \to \gamma$ добавляют $(B \to \bullet \gamma;\,[j,j])$ в $\mathrm{chart}[j]$, где $j$ --- конец уже разобранного префикса состояния.

	\item[\textsc{Scanner}] применяют, если справа от $\bullet$ --- терминал/часть речи $B$ и $w_i$ согласуется с $B$.
	Создают $(B \to w_i \bullet;\,[i,i{+}1])$ в $\mathrm{chart}[i{+}1]$.

	\item[\textsc{Completer}] применяют, если $\bullet$ в конце правила: $(B \to \gamma \bullet;\,[j,k])$ завершено.
	Для каждого $(A \to \alpha \bullet B \beta;\,[i,j])$ в $\mathrm{chart}[j]$ добавляют $(A \to \alpha B \bullet \beta;\,[i,k])$ в $\mathrm{chart}[k]$.
\end{description}

Дубликаты состояний не добавляют (\textsc{Enqueue} только при отсутствии состояния в множестве).

\subsubsection{Результат}

Разбор успешен, если в $\mathrm{chart}[n]$ есть состояние $(\gamma \to S \bullet;\,[0,n])$ (или эквивалентно завершённое $S$, покрывающее всё предложение).
Чарт неявно кодирует \emph{все} синтаксические деревья; поддеревья хранятся один раз и переиспользуются.

Для извлечения конкретного дерева модифицируют \textsc{Completer}, сохраняя ссылки на породившие состояния (аналогично обратным указателям в CKY).

\subsubsection{Сложность и сравнение с CKY}

\begin{itemize}
	\item общий случай: $O(n^3)$ по длине $n$;
	\item для неоднозначных грамматик на практике часто ближе к $O(n^2)$;
	\item для некоторых классов грамматик (например, $LR(k)$) --- линейное время;
	\item в отличие от CKY, не нужна НФХ; зато реализация и анализ состояний сложнее.
\end{itemize}

В~\cite{jurafsky2026speech}, гл.~13, CKY излагается как основной табличный bottom-up алгоритм; алгоритм Эрли и chart parsing упоминаются как родственные подходы на основе того же принципа мемоизации поддеревьев.

% Источники: report.pdf (разд.~2.3.1), Jurafsky \& Martin, 3rd draft, гл.~13;
% Earley — Earley (1970); в 3-м черновике Jurafsky (2026) — краткое упоминание в гл.~13.

\subsection{Алгоритм Кока--Янгера--Касами (CKY)}

\subsubsection{Что можно сохранить из текущего описания (\texttt{report.pdf}, разд.~2.3.1)}

Следующие положения согласуются с изложением в~\cite{jurafsky2026speech}, гл.~13, и могут быть оставлены (после небольшой правки формулировок):
\begin{itemize}
	\item алгоритм строит \emph{дерево составляющих} (конституентный разбор) для предложений, порождаемых контекстно-свободной грамматикой;
	\item перед разбором грамматику приводят к \emph{нормальной форме Хомского} (НФХ): правила вида $A \to BC$ или $A \to w$~\cite{jurafsky2026speech};
	\item используется \emph{динамическое программирование} на верхнетреугольной таблице размера $(n{+}1)\times(n{+}1)$ для предложения из $n$ слов;
	\item ответ разбора (все допустимые корневые конституенты) содержится в ячейке $[0,n]$;
	\item для \emph{восстановления дерева} целесообразно хранить указатели на подынтервалы $(i,k)$ и $(k,j)$, породившие данную запись;
	\item \emph{временная сложность} распознавания --- $O(n^3)$ при фиксированной грамматике (точнее $O(n^3|G|)$, если учитывать размер грамматики).
\end{itemize}

\subsubsection{Замечания и ошибки относительно~\cite{jurafsky2026speech}, гл.~13}

\begin{enumerate}
	\item \textbf{Название алгоритма.}
	В отчёте: «Кок--Янгер--Касами--Шварц--Зиппель».
	В гл.~13 Jurafsky --- классический \textbf{CKY / CYK} (Cocke--Kasami--Younger)~\cite{younger1967recognition}.
	Расширение Schwartz--Zippel --- отдельное современное направление~\cite{makarov2022cyk}; для базового парсера его включать в название не следует.

	\item \textbf{Рекуррентное соотношение (существенная ошибка).}
	В отчёте:
	\[
		dp_{i,j} = dp_{i,j} \cup \bigl(dp_{i,k} \cup dp_{k,j}\bigr), \quad i<k<j,
	\]
	что не задаёт порождение новых нетерминалов.
	По~\cite{jurafsky2026speech}, \S13.2.2--13.2.3, ячейка $[i,j]$ --- это \emph{множество нетерминалов}, покрывающих слова с позиции $i$ по $j$.
	Для каждого разбиения $i<k<j$ и каждого правила $A \to B\,C$ из НФХ, если $B \in \mathrm{cell}[i,k]$ и $C \in \mathrm{cell}[k,j]$, добавляют $A$ в $\mathrm{cell}[i,j]$.

	\item \textbf{Семантика ячеек.}
	В отчёте смешаны «тип объединения» и объединение множеств без ссылки на правила грамматики.
	В Jurafsky объединяются только те нетерминалы, которые \emph{разрешены грамматикой} при бинарном разбиении интервала.

	\item \textbf{Порядок заполнения таблицы.}
	Формулировка «слева направо, снизу вверх» неточна.
	В гл.~13: внешний цикл --- по \emph{столбцам} (длине интервала) слева направо, внутри столбца --- по \emph{строкам} снизу вверх; для ячейки $[i,j]$ перебирают все точки разбиения $k$ ($i<k<j$).

	\item \textbf{Индексация.}
	У Jurafsky индексы $0,\ldots,n$ --- «столбики» \emph{между} словами; слова дают заполнение диагонали вида $[i,i{+}1]$.
	В отчёте $i,j = \overline{0,n}$ при $i<j$ без явного различения ячеек длины~1 и произвольного span --- это стоит уточнить.

	\item \textbf{Распознавание и разбор.}
	Базовый CKY --- \emph{распознаватель}: успех, если $S \in \mathrm{cell}[0,n]$ (\S13.2.3).
	Фраза «варианты разбора в $dp_{0,n}$» верна только если в ячейках хранят \emph{все} нетерминалы и (для полного парсера) обратные указатели; иначе это лишь множество возможных корней.

	\item \textbf{Дисамбигуация.}
	«Вручную или предобученным классификатором» --- не часть классического CKY в \S13.2.
	Выбор лучшего разбора Jurafsky выносит в \S13.3+ (нейросетевой span-based parsing, вероятностный CKY); для диплома формулировку лучше отделить от описания ядра алгоритма.

	\item \textbf{Ссылки на доказательство сложности.}
	Доказательство $O(n^3)$ для CYK --- Younger (1967), Kasami (1965); в~\cite{makarov2022cyk} --- обзор родственных алгоритмов, а не классическое доказательство.
	Теорему о приведении к НФХ логичнее ссылать на гл.~12~\cite{jurafsky2026speech}, а не только на Makarov.

	\item \textbf{Опечатки в соседнем тексте.}
	«завсисмостей» (табл.~сравнения), «было было» (восстановление дерева).
\end{enumerate}

\subsubsection{Исправленное описание (по~\cite{jurafsky2026speech}, гл.~13)}

Пусть вход --- слова $w_1,\ldots,w_n$, грамматика $G=(N,\Sigma,P,S)$ в НФХ.
Таблица $\mathrm{cell}[i,j]$ ($0 \le i < j \le n$) --- множество нетерминалов, порождающих $w_{i+1}\ldots w_j$.

\paragraph{Инициализация.}
Для каждого $i=0,\ldots,n{-}1$:
$\mathrm{cell}[i,i{+}1] \gets \{A \mid (A \to w_{i+1}) \in P\}$ (и при необходимости частеречные категории из лексикона).

\paragraph{Заполнение.}
Для $j = 2,\ldots,n$ (длина интервала), для $i = j{-}2,\ldots,0$ (снизу вверх в столбце $j$), для $k = i{+}1,\ldots,j{-}1$:
\[
	\mathrm{cell}[i,j] \gets \mathrm{cell}[i,j] \cup
	\{A \mid A \to B\,C \in P,\; B \in \mathrm{cell}[i,k],\; C \in \mathrm{cell}[k,j]\}.
\]

\paragraph{Распознавание.}
Предложение принимается грамматикой тогда и только тогда, когда $S \in \mathrm{cell}[0,n]$.

\paragraph{Построение дерева.}
При добавлении $A$ из правила $A \to B\,C$ с разбиением $k$ сохраняют указатель $(B,C,k)$.
Восстановление --- рекурсивный обход от $S \in \mathrm{cell}[0,n]$.

\paragraph{Сложность.}
При $|G|$ правил и длине $n$ --- $O(n^3|G|)$; при фиксированной грамматике --- $O(n^3)$.

%--------------------------------------------------------------------
\subsection{Алгоритм Эрли (Earley)}

% В 3-м черновике Jurafsky (2026) алгоритм Эрли упоминается в гл.~13
% как родственный подход; детали ниже по~\cite{earley1970efficient}.

Алгоритм Эрли~\cite{earley1970efficient} --- алгоритм \emph{динамического программирования} для разбора произвольных КС-грамматик (НФХ \emph{не требуется}).
Он совмещает нисходящие ожидания и восходящее подтверждение конституентов; в отличие от CKY, заполнение идёт \emph{слева направо} по позициям входа.

\subsubsection{Структуры данных}

Для предложения длины $n$ строят \textbf{чарт} (chart) --- массив из $n{+}1$ множеств состояний $\mathrm{chart}[0],\ldots,\mathrm{chart}[n]$.
Позиция $i$ соответствует \emph{зазору} перед $i$-м словом (между словами).

\textbf{Состояние} (dotted rule) записывают как тройку
\[
	(A \to \alpha \bullet \beta;\; [i,j]),
\]
где $A \to \alpha\beta$ --- правило грамматики, $\bullet$ --- позиция в правой части, $[i,j]$ --- интервал входа, который уже покрыт префиксом $\alpha$ (от $i$ до текущей позиции обработки).

Стартовое состояние: $(\gamma \to \bullet S;\,[0,0])$, где $\gamma$ --- фиктивный начальный символ.

\subsubsection{Операторы}

Для каждой позиции $i$ обрабатывают состояния в $\mathrm{chart}[i]$ в порядке добавления (без удаления и без возврата назад).

\begin{description}
	\item[\textsc{Predictor}] применяют, если справа от $\bullet$ стоит нетерминал $B$ (не часть речи).
	Для каждого правила $B \to \gamma$ добавляют $(B \to \bullet \gamma;\,[j,j])$ в $\mathrm{chart}[j]$, где $j$ --- конец уже разобранного префикса состояния.

	\item[\textsc{Scanner}] применяют, если справа от $\bullet$ --- терминал/часть речи $B$ и $w_i$ согласуется с $B$.
	Создают $(B \to w_i \bullet;\,[i,i{+}1])$ в $\mathrm{chart}[i{+}1]$.

	\item[\textsc{Completer}] применяют, если $\bullet$ в конце правила: $(B \to \gamma \bullet;\,[j,k])$ завершено.
	Для каждого $(A \to \alpha \bullet B \beta;\,[i,j])$ в $\mathrm{chart}[j]$ добавляют $(A \to \alpha B \bullet \beta;\,[i,k])$ в $\mathrm{chart}[k]$.
\end{description}

Дубликаты состояний не добавляют (\textsc{Enqueue} только при отсутствии состояния в множестве).

\subsubsection{Результат}

Разбор успешен, если в $\mathrm{chart}[n]$ есть состояние $(\gamma \to S \bullet;\,[0,n])$ (или эквивалентно завершённое $S$, покрывающее всё предложение).
Чарт неявно кодирует \emph{все} синтаксические деревья; поддеревья хранятся один раз и переиспользуются.

Для извлечения конкретного дерева модифицируют \textsc{Completer}, сохраняя ссылки на породившие состояния (аналогично обратным указателям в CKY).

\subsubsection{Сложность и сравнение с CKY}

\begin{itemize}
	\item общий случай: $O(n^3)$ по длине $n$;
	\item для неоднозначных грамматик на практике часто ближе к $O(n^2)$;
	\item для некоторых классов грамматик (например, $LR(k)$) --- линейное время;
	\item в отличие от CKY, не нужна НФХ; зато реализация и анализ состояний сложнее.
\end{itemize}

В~\cite{jurafsky2026speech}, гл.~13, CKY излагается как основной табличный bottom-up алгоритм; алгоритм Эрли и chart parsing упоминаются как родственные подходы на основе того же принципа мемоизации поддеревьев.
